%---------------------------------------------------------------
\chapter{\kesimpulan}
%---------------------------------------------------------------

%---------------------------------------------------------------
\section{Kesimpulan}
%---------------------------------------------------------------

Berdasarkan hasil perancangan, implementasi, dan pengujian yang telah dilakukan pada pengembangan website ini, dapat diambil beberapa kesimpulan sebagai berikut:

\begin{enumerate}
	\item Website berhasil dirancang dan diimplementasikan sesuai dengan kebutuhan yang telah didefinisikan pada tahap analisis kebutuhan sistem. Website ini digunakan sebagai \textit{landing page} atau media perkenalan Laboratorium IR NLP UI. 
	
	\item Sistem yang dibangun mampu menjalankan fitur-fitur utama, seperti manajemen data, pencarian informasi, dan integrasi API dengan baik serta sesuai dengan tujuan awal pengembangan.
	
	\item Berdasarkan hasil pengujian fungsional dan non-fungsional, sistem menunjukkan performa yang stabil serta dapat diakses dengan baik melalui berbagai perangkat dan peramban (\textit{browser}).
	
	\item Penggunaan teknologi seperti Django, HTML, CSS, JavaScript, PostgreSQL, API OpenAlex, dan APScheduler terbukti mendukung pengembangan sistem yang terstruktur, modular, dan mudah dikembangkan kembali (scalable).
	
	\item Website yang dikembangkan mampu membantu mengenalkan Laboratorium IR NLP UI kepada khalayak luas serta menyediakan informasi yang relevan dan terkini mengenai kegiatan dan publikasi laboratorium.
\end{enumerate}


%---------------------------------------------------------------
\section{Saran}
%---------------------------------------------------------------

Untuk pengembangan dan penyempurnaan website di masa mendatang, terdapat beberapa saran yang dapat dipertimbangkan sebagai berikut:

\begin{enumerate}
	\item Menambahkan fitur keamanan lanjutan, seperti enkripsi data, perlindungan terhadap serangan \textit{brute-force}, serta penerapan \textit{two-factor authentication} (2FA) guna meningkatkan keamanan sistem.
	
	\item Mengoptimalkan performa website, khususnya dalam hal waktu pemuatan halaman (\textit{loading time}) dan pengelolaan \textit{query} basis data agar lebih efisien ketika jumlah pengguna meningkat.
	
	\item Menambahkan fitur pelaporan (\textit{reporting}) dan visualisasi data dalam bentuk grafik atau \textit{dashboard} interaktif untuk memudahkan analisis data.
	
	\item Menambahkan demo interaktif terkait IR dan NLP yang dapat digunakan oleh pengunjung untuk memahami konsep-konsep dasar dalam bidang tersebut.
	
	\item Melakukan pengujian berkelanjutan (\textit{continuous testing}) serta pemeliharaan sistem secara berkala agar website tetap berjalan optimal dan terhindar dari kesalahan (\textit{bug}) di masa mendatang.
\end{enumerate}
